% 学校模板的10类标识和8种著录示意，集中由第一章载入。
% 字段占位不是实际文献，不进入 BibTeX 列表；不受旧备注开关控制。
\subsection{十种文献类型标识}

学校模板列示的类型为：期刊文章[J]，专著[M]，论文集[C]，学位论文[D]，报告[R]，标准[S]，专利[P]，网上文章[EB/OL]，报纸文章[N]，论文集中的析出文献[A]。整本论文集与其中的析出文献应按实际来源选择类型。

\subsection{八种著录格式示意}

以下保留学校模板列示的八种著录格式。“序号”“作者”“题名”等均是字段占位，不是真实文献信息，不加入实际参考文献列表，也不参与其编号。报告和析出文献属于上述类型清单；它们未被另行计作学校原表的八种示例。

\noindent\textbf{期刊文章：} [序号] 作者. 文章名[J]. 刊名（外文期刊可缩写），年，卷号（期号）：起～止页码。

\noindent\textbf{专（译）著：} [序号] 作者. 书名[M]. （译者，译）. 出版地：出版者，出版年. 起～止页码。

\noindent\textbf{论文集：} [序号] 作者. 题名. 论文集名[C]. 出版地：出版者，出版年. 起～止页码。

\noindent\textbf{学位论文：} [序号] 姓名. 论文题目[D]. 授予地：授予单位，授予年。

\noindent\textbf{专利：} [序号] 申请者. 专利名：国名，专利号[P]. 出版日期。

\noindent\textbf{网上文章：} [序号] 作者，文章名[EB/OL]. 文章出处或可获得网址，发表或更新日期。

\noindent\textbf{技术标准：} [序号] 发布单位. 技术标准编号. 技术标准名称[S]. 出版地：出版者，出版日期。

\noindent\textbf{报刊文章：} [序号] 作者. 文章名[N]. 报刊名，出版日期（版次）。
